\documentclass[a4paper,12pt]{article}
\usepackage[T1]{fontenc}
\usepackage[utf8]{inputenc}
\usepackage[czech]{babel}
\usepackage{amsmath}
\usepackage{enumitem}
\usepackage{hyperref}
\usepackage{geometry}
\geometry{margin=2.5cm}
\setlist{itemsep=2pt,topsep=4pt}

\title{Dokumentace k projektu IFJ/IAL}
\author{xpocarj00 \\ xuherto00 \\ xchalo21 \\ xkravc03}
\date{Listopad 2025}

\begin{document}

\maketitle

Dokumentace k projektu z předmětu IFJ a IAL\\
Implementace překladače pro jazyk IFJ24\\
Tým xpocarj00, xuherto00, xchalo21, xkravc03, varianta TRP-izp\\
Každý člen přispěl 25\% práce.

\tableofcontents
\newpage

\section{Rozdělení práce mezi členy týmu}

\subsection{xpocarj00 (25\%)}
\begin{itemize}
    \item Lexikální analýza
    \item Generátor kódu
\end{itemize}

\subsection{xuherto00 (25\%)}
\begin{itemize}
    \item Syntaktická analýza
    \item Precedenční logika
    \item Tabulka symbolů
\end{itemize}

\subsection{xchalo21 (25\%)}
\begin{itemize}
    \item Tabulka symbolů
    \item Generátor kódu
\end{itemize}

\subsection{xkravc03 (25\%)}
\begin{itemize}
    \item Testovací soubory
    \item Makefile
    \item Dokumentace
    \item Pomocné struktury seznam a zásobník
\end{itemize}

\section{Způsob řešení jednotlivých částí projektu}

\subsection{Lexikální analýza}
Lexikální analýza je začátkem a základem překladače. Není nejkomplikovanější, ale jakákoli chyba rychle vede k selhání celého programu, který obdrží nesprávná data. V našem projektu ji realizují soubory \texttt{scanner.c} a \texttt{scanner.h}. Návrh kódu vychází z konečného automatu, který průběžně doplňujeme o nově zjištěné požadavky. Pomocí \texttt{getchar()} načítáme symbol, jenž zpracujeme podle aktuálního stavu; podle potřeby přepínáme stavy a zapisujeme symbol do tokenu. K \texttt{scanner.c} přistupujeme funkcí \texttt{Token* getToken()}, která vrací odkaz na nově vytvořený token obsahující typ a atribut reprezentovaný ukazatelem na pole znaků (maximální velikost \texttt{TOKEN\_STRING\_SIZE}), číslem typu \texttt{int}, \texttt{float}, nebo hodnotou typu \texttt{enum}.

\subsection{Syntaktická analýza}
Nejdůležitější částí programu je syntaktická analýza implementovaná pomocí syntaktického analyzátoru v souborech \texttt{parser.c} a \texttt{parser.h}. Samotný analyzátor spouští funkce \texttt{NProg}. Tato funkce je řízena LL gramatikou a metodou rekurzivního sestupu přes LL tabulku. Postupně prochází tokeny uložené ve strukturovaném seznamu \texttt{TList} definovaném v souborech \texttt{list.c} a \texttt{list.h}. Vyhodnocuje vnořené funkce sestavené pro každý neterminál, které vracejí hodnotu \texttt{true}/\texttt{false} podle korektnosti vstupního programu. Návratové hodnoty těchto funkcí se propagují zpět do \texttt{NProg}, která vrátí \texttt{true}, pokud je vstup syntakticky správný, nebo \texttt{false} v opačném případě.

\subsection{Zpracování výrazů pomocí precedenční logiky}
Precedenční logiku má na starosti funkce \texttt{precProces} v souborech \texttt{precedencna.c} a \texttt{precedencna.h}, která se volá pro každý token ve výrazu. Na strukturovaný zásobník \texttt{StackNode} přidá token a může provést zjednodušení podle pravidel precedenční tabulky. Použitá pravidla ukládá na samostatný zásobník, jenž je následně předán tokenu před výrazem pro jednodušší přístup a manipulaci. Po celý čas probíhá typová kontrola, aby byla zajištěna správnost výsledné hodnoty.

\subsection{Generátor kódu}
Generování kódu zajišťuje funkce \texttt{codeGen} v souborech \texttt{interpret.c} a \texttt{interpret.h}. Tato funkce prochází seznam tokenů podle předdefinované logiky. Podle prvního tokenu volá specializované funkce pro daný příkaz. Každá z nich využívá pomocné tiskové funkce, které na \texttt{stdout} generují trojadresné instrukce pro interpret \texttt{IFJcode24}. Současně využívá strukturu \texttt{StackNode}, do níž ukládá pořadí zanoření příkazů, například \texttt{while} a \texttt{if}. Při práci s výrazy se používají pravidla precedenční logiky uložená v položce \texttt{rules} předchozího tokenu, čímž se zaručí správný výpočet.

\subsection{Tabulka symbolů TRP-izp}
Tabulka symbolů je implementována jako strukturovaná hashovací tabulka \texttt{HashItem} podle varianty zadání TRP-izp v souborech \texttt{symtable.c} a \texttt{symtable.h}. Má na starosti spravovat uživatelem definované proměnné a funkce spolu s předdefinovanými funkcemi \texttt{ifj.xxx}. Tabulka používá "magic number" \texttt{hashNumber} 5381 pro optimální mapování položek. Obsahuje informace o typu \texttt{const}/\texttt{var}, parametrech funkcí ve struktuře \texttt{symItems}, lokální skupině, úrovni zanoření a o tom, zda byla proměnná inicializována nebo definována.

\subsection{LL gramatika}
Aktuálně používaná LL gramatika je uvedena níže. V této verzi dokumentu uvádíme pouze textový zápis pravidel.

\begin{enumerate}[label=\arabic*.]
    \item $\langle NPROG \rangle \rightarrow \langle NPROLOG \rangle \langle NDEF \rangle EOF$
    \item $\langle NPROLOG \rangle \rightarrow \text{const } ID = @inport(\text{\textquotedblleft}ifj24.zig\text{\textquotedblright});$
    \item $\langle NDEF \rangle \rightarrow \text{pub fn } ID (\langle NDEFPARAM \rangle) \langle NTYPEFUN \rangle \{\langle NNEXT \rangle\} \langle NDEF \rangle$
    \item $\langle NDEF \rangle \rightarrow \varepsilon$
    \item $\langle NDEFPARAM \rangle \rightarrow ID : \langle NTYPE \rangle \langle NDEFCOMMA \rangle$
    \item $\langle NDEFPARAM \rangle \rightarrow \varepsilon$
    \item $\langle NTYPE \rangle \rightarrow I32$
    \item $\langle NTYPE \rangle \rightarrow F64$
    \item $\langle NTYPE \rangle \rightarrow []u8$
    \item $\langle NTYPE \rangle \rightarrow ?I32$
    \item $\langle NTYPE \rangle \rightarrow ?F64$
    \item $\langle NTYPE \rangle \rightarrow ?[]u8$
    \item $\langle NTYPEFUN \rangle \rightarrow \langle NTYPE \rangle$
    \item $\langle NTYPEFUN \rangle \rightarrow void$
    \item $\langle NDEFCOMMA \rangle \rightarrow ,\langle NDEFPARAM \rangle$
    \item $\langle NDEFCOMMA \rangle \rightarrow \varepsilon$
    \item $\langle NNEXT \rangle \rightarrow \langle NCOMMAND \rangle$
    \item $\langle NNEXT \rangle \rightarrow \varepsilon$
    \item $\langle NCOMMAND \rangle \rightarrow if(\langle NEXPRESS \rangle)\langle NIDNULL \rangle\{\langle NNEXT \rangle\}\langle NELSE \rangle\langle NNEXT \rangle$
    \item $\langle NCOMMAND \rangle \rightarrow while(\langle NEXPRESS \rangle)\langle NIDNULL \rangle\{\langle NNEXT \rangle\}\langle NNEXT \rangle$
    \item $\langle NCOMMAND \rangle \rightarrow return\langle NRETVALUE \rangle;$
    \item $\langle NCOMMAND \rangle \rightarrow \langle NVAR \rangle\langle NNEXT \rangle$
    \item $\langle NCOMMAND \rangle \rightarrow ID\langle NSUPPORT \rangle;\langle NNEXT \rangle$
    \item $\langle NRETVALUE \rangle \rightarrow \langle NVALUE \rangle$
    \item $\langle NRETVALUE \rangle \rightarrow \varepsilon$
    \item $\langle NSUPPORT \rangle \rightarrow (\langle NPARAM \rangle)$
    \item $\langle NSUPPORT \rangle \rightarrow =\langle NVALUE \rangle$
    \item $\langle NVAR \rangle \rightarrow \text{const } ID : \langle NTYPEVAR \rangle = \langle NVALUE \rangle;$
    \item $\langle NVAR \rangle \rightarrow \text{var } ID : \langle NTYPEVAR \rangle = \langle NVALUE \rangle;$
    \item $\langle NTYPEVAR \rangle \rightarrow :\langle NTYPE \rangle$
    \item $\langle NTYPEVAR \rangle \rightarrow \varepsilon$
    \item $\langle NVALUE \rangle \rightarrow \langle NVALUEX \rangle\langle NSIGN \rangle\langle NVALUEX \rangle$
    \item $\langle NVALUE \rangle \rightarrow \langle NVALUEX \rangle$
    \item $\langle NVALUE \rangle \rightarrow null$
    \item $\langle NVALUEX \rangle \rightarrow (\langle NVALUE \rangle)$
    \item $\langle NVALUEX \rangle \rightarrow ID$
    \item $\langle NVALUEX \rangle \rightarrow ID(\langle NPARAM \rangle)$
    \item $\langle NPARAM \rangle \rightarrow ID\langle NCOMMA \rangle$
    \item $\langle NPARAM \rangle \rightarrow \varepsilon$
    \item $\langle NCOMMA \rangle \rightarrow ,\langle NPARAM \rangle$
    \item $\langle NCOMMA \rangle \rightarrow \varepsilon$
    \item $\langle NEXPRESS \rangle \rightarrow ID$
    \item $\langle NEXPRESS \rangle \rightarrow \langle NVALUE \rangle\langle NCOMPSIGN \rangle\langle NVALUE \rangle$
    \item $\langle NIDNULL \rangle \rightarrow |ID|$
    \item $\langle NSIGN \rangle \rightarrow +$
    \item $\langle NSIGN \rangle \rightarrow -$
    \item $\langle NSIGN \rangle \rightarrow *$
    \item $\langle NSIGN \rangle \rightarrow /$
    \item $\langle NCOMPSIGN \rangle \rightarrow ==$
    \item $\langle NCOMPSIGN \rangle \rightarrow !=$
    \item $\langle NCOMPSIGN \rangle \rightarrow >$
    \item $\langle NCOMPSIGN \rangle \rightarrow >=$
    \item $\langle NCOMPSIGN \rangle \rightarrow <=$
    \item $\langle NCOMPSIGN \rangle \rightarrow <$
    \item $\langle NELSE \rangle \rightarrow else\{\langle NNEXT \rangle\}$
    \item $\langle NELSE \rangle \rightarrow \varepsilon$
\end{enumerate}

\subsection{LL tabulka}
Konkrétní LL tabulka bude doplněna v další verzi dokumentace. V současnosti shromažďujeme výsledky validace, aby bylo možné tabulku prezentovat v přesné a finální podobě.

\subsection{Precedenční tabulka}
Stejně jako u LL tabulky bude kompletní precedenční tabulka zařazena po finalizaci verifikace všech pravidel. Textová podoba pravidel je již implementována v kódu a bude zpracována do přehledné tabulky.

\subsection{Konečný automat}
Grafické znázornění konečného automatu připravujeme pro budoucí revizi dokumentace. V této verzi poskytujeme pouze textový popis zpracování tokenů, aby dokument zůstal konzistentní.

\end{document}
